\section{Constrained finite time optimal control}

In the previous chapter we studied finite- and infinite-horizon linear quadratic control in the
absence of inequality constraints. That setting already showed the predictive-control logic very
clearly: one optimizes over a future input sequence, predicts the resulting state trajectory, and
uses the optimizer to determine the control action. However, one of the main reasons why MPC
is so useful in practice is precisely that it can handle constraints explicitly. Real systems cannot
operate with arbitrarily large inputs, and states or outputs are often restricted by safety,
performance, or physical limitations.\\

The ideal constrained problem would be an infinite-horizon optimal control problem of the form
\[
J_\infty^\star(x(0))
=
\min_{u(\cdot)}
\sum_{i=0}^{\infty} \stagecost(x_i,u_i)
\]
subject to
\[
x_{i+1}=Ax_i+Bu_i,\qquad i=0,1,\dots,
\]
together with the constraints
\[
x_i\in\cX,\qquad u_i\in\cU,\qquad i=0,1,\dots,
\]
and the initial condition
\[
x_0=x(0).
\]
This problem is conceptually attractive because it captures the long-term consequences of the
control action while respecting the constraints at all times. The difficulty is that it involves an
infinite number of decision variables, so it is generally not directly computable online.\\

The standard MPC idea is therefore to replace this ideal problem by a constrained finite time
optimal control problem over a horizon of length \(N\). At time \(k\), one solves
\[
J^\star(x(k))
=
\min_{U}
\terminalcost(x_N)+\sum_{i=0}^{N-1}\stagecost(x_i,u_i)
\]
subject to
\[
x_{i+1}=Ax_i+Bu_i,\qquad i=0,\dots,N-1,
\]
\[
x_i\in\cX,\qquad u_i\in\cU,\qquad i=0,\dots,N-1,
\]
\[
x_N\in \cX_f,
\]
\[
x_0=x(k),
\]
where
\[
U=\{u_0,\dots,u_{N-1}\}.
\]
The finite horizon makes the problem computationally tractable, while the terminal cost
\(\terminalcost(x_N)\) and terminal set \(\cX_f\) are introduced to approximate the effect of the
tail beyond the prediction horizon.

\subsection{General CFTOC formulation}

Throughout this chapter we consider a discrete-time linear system
\[
x_{i+1}=Ax_i+Bu_i,
\]
with state and input constrained to lie in polyhedral sets
\[
\cX=\{x\in\R^n\mid A_xx\le b_x\},
\qquad
\cU=\{u\in\R^m\mid A_uu\le b_u\},
\]
and terminal set
\[
\cX_f=\{x\in\R^n\mid A_fx\le b_f\}.
\]
The associated cost has the general form
\[
J(x_0,U)=\terminalcost(x_N)+\sum_{i=0}^{N-1}\stagecost(x_i,u_i).
\]
Two important classes of cost functions will be considered. In the quadratic case, one uses
\[
\terminalcost(x_N)=x_N^\top P x_N,
\qquad
\stagecost(x_i,u_i)=x_i^\top Q x_i+u_i^\top R u_i,
\]
with
\[
P\succeq 0,\qquad Q\succeq 0,\qquad R\succ 0.
\]
In the norm-based case, one uses either the \(1\)-norm or the \(\infty\)-norm:
\[
\terminalcost(x_N)=\norm{P x_N}_p,
\qquad
\stagecost(x_i,u_i)=\norm{Qx_i}_p+\norm{Ru_i}_p,
\qquad p\in\{1,\infty\},
\]
where the weighting matrices are chosen with appropriate rank assumptions so that the cost
meaningfully penalizes the variables.

\subsection{Feasible set}

An important object associated with the constrained finite-horizon problem is its feasible set.
It describes from which initial states the optimization problem can actually be solved while
respecting all constraints over the horizon. For a given horizon \(N\), the feasible set is
\[
\cX_N
=
\Bigl\{
x_0\in\R^n\;\Big|\;
\exists (u_0,\dots,u_{N-1}) \text{ such that }
x_i\in\cX,\;
u_i\in\cU,\;
i=0,\dots,N-1,
\;
x_N\in\cX_f,
\;
x_{i+1}=Ax_i+Bu_i
\Bigr\}.
\]
This set is fundamental because outside \(\cX_N\) the CFTOC problem is simply infeasible.
Notice also that \(\cX_N\) depends on the dynamics, the constraints, the terminal set, and the
horizon length, but not on the specific choice of cost. The cost only selects, among the feasible
trajectories, the one that is considered optimal.\\
This is in clear contrast with the unconstrained LQR setting. There, every initial state is
admissible because there are no state or input limits, and the solution has a linear feedback form.
Once constraints are introduced, this global linear structure is lost, and the control law becomes
more intricate.

\section{Quadratic-cost CFTOC and quadratic programming}

We first consider the constrained finite-horizon problem with quadratic stage and terminal cost:
\[
J^\star(x(k))
=
\min_{U}
x_N^\top P x_N + \sum_{i=0}^{N-1}\left(x_i^\top Q x_i + u_i^\top R u_i\right)
\]
subject to
\[
x_{i+1}=Ax_i+Bu_i,\qquad i=0,\dots,N-1,
\]
\[
x_i\in\cX,\qquad u_i\in\cU,\qquad i=0,\dots,N-1,
\]
\[
x_N\in\cX_f,
\qquad
x_0=x(k).
\]
Because the dynamics are linear, the constraints are polyhedral, and the cost is quadratic, this
problem can be rewritten as a quadratic program. This is one of the most important structural
facts in nominal linear MPC.

\subsection{QP formulation without eliminating the states}

A first way to rewrite the quadratic-cost CFTOC problem as a quadratic program is to
\emph{keep the predicted states as decision variables} and impose the system dynamics as
explicit equality constraints. This formulation is often attractive in practice because it preserves
the sparse structure of the dynamics.\\
Starting from
\[
J^\star(x(k))
=
\min_U
x_N^\top P x_N + \sum_{i=0}^{N-1}\left(x_i^\top Q x_i + u_i^\top R u_i\right)
\]
subject to
\[
x_{i+1}=Ax_i+Bu_i,\qquad i=0,\dots,N-1,
\]
\[
x_i\in\cX,\qquad u_i\in\cU,\qquad i=0,\dots,N-1,
\]
\[
x_N\in\cX_f,\qquad x_0=x(k),
\]
we introduce the stacked decision vector
\[
z=
\begin{bmatrix}
x_1^\top & \dots & x_N^\top & u_0^\top & \dots & u_{N-1}^\top
\end{bmatrix}^\top .
\]
With this choice of variables, the CFTOC problem can be transformed into a QP of the form
\[
J^\star(x(k))
=
\min_z
\begin{bmatrix}
z^\top & x(k)^\top
\end{bmatrix}
\begin{bmatrix}
\bar H & 0\\
0 & Q
\end{bmatrix}
\begin{bmatrix}
z\\
x(k)
\end{bmatrix}
\]
subject to
\[
G_{\mathrm{in}} z \le w_{\mathrm{in}} + E_{\mathrm{in}}x(k),
\qquad
G_{\mathrm{eq}} z = E_{\mathrm{eq}}x(k).
\]
The equality constraints come directly from the dynamics
\[
x_{i+1}=Ax_i+Bu_i.
\]
Collecting them over the whole horizon gives
\[
G_{\mathrm{eq}} z = E_{\mathrm{eq}}x(k),
\]
with the block structure
\[
G_{\mathrm{eq}}=
\begin{bmatrix}
I      & 0      & \cdots & 0      & -B     & 0      & \cdots & 0\\
-A     & I      & \cdots & 0      & 0      & -B     & \cdots & 0\\
0      & -A     & \ddots & \vdots & \vdots & \ddots & \ddots & \vdots\\
\vdots & \ddots & \ddots & I      & 0      & \cdots & 0      & -B
\end{bmatrix},
\qquad
E_{\mathrm{eq}}=
\begin{bmatrix}
A\\
0\\
\vdots\\
0
\end{bmatrix}.
\]
The first block row corresponds to
\[
x_1-Ax(k)-Bu_0=0,
\]
while the remaining rows encode
\[
x_{i+1}-Ax_i-Bu_i=0,\qquad i=1,\dots,N-1.
\]
Next, suppose the constraint sets are given in polyhedral form as
\[
\cX=\{x\mid A_xx\le b_x\},\qquad
\cU=\{u\mid A_uu\le b_u\},\qquad
\cX_f=\{x\mid A_fx\le b_f\}.
\]
Then all state, input, and terminal constraints can be collected into
\[
G_{\mathrm{in}} z \le w_{\mathrm{in}} + E_{\mathrm{in}}x(k).
\]
The matrix \(G_{\mathrm{in}}\) has a block-diagonal structure, with the state constraints acting
on \(x_1,\dots,x_{N-1}\), the terminal constraint acting on \(x_N\), and the input constraints
acting on \(u_0,\dots,u_{N-1}\). Explicitly,
\[
G_{\mathrm{in}}=
\begin{bmatrix}
0 & 0 & \cdots & 0 & 0 & 0 & \cdots & 0\\
A_x & 0   & \cdots & 0   & 0   & 0   & \cdots & 0\\
0   & A_x & \cdots & 0   & 0   & 0   & \cdots & 0\\
\vdots & \vdots & \ddots & \vdots & \vdots & \vdots & \ddots & \vdots\\
0   & 0   & \cdots & A_x & 0   & 0   & \cdots & 0\\
0   & 0   & \cdots & A_f & 0   & 0   & \cdots & 0\\[0.3em]
0   & 0   & \cdots & 0   & A_u & 0   & \cdots & 0\\
0   & 0   & \cdots & 0   & 0   & A_u & \cdots & 0\\
\vdots & \vdots &  & \vdots & \vdots & \vdots & \ddots & \vdots\\
0   & 0   & \cdots & 0   & 0   & 0   & \cdots & A_u
\end{bmatrix},
\]
\[
w_{\mathrm{in}}=
\begin{bmatrix}
b_x \\
b_x\\
b_x\\
\vdots\\
b_x\\
b_f\\
b_u\\
b_u\\
\vdots\\
b_u
\end{bmatrix},
\qquad
E_{\mathrm{in}}=
\begin{bmatrix}
-A_x\\
0\\
\vdots\\
0
\end{bmatrix}.
\]
The term \(E_{\mathrm{in}}x(k)\) appears because the initial state constraint \(x_0\in\cX\) is not
part of the decision vector \(z\), so its contribution must be written separately on the
right-hand side.\\
Finally, the cost is obtained directly from
\[
x_N^\top P x_N+\sum_{i=0}^{N-1}\left(x_i^\top Qx_i+u_i^\top Ru_i\right).
\]
Since \(x_0=x(k)\) is fixed and not optimized over, the quadratic term in the optimization
variables is described by the block-diagonal matrix
\[
\bar H=
\mathrm{blockdiag}(Q,\dots,Q,P,R,\dots,R),
\]
that is,
\[
\bar H=
\begin{bmatrix}
Q &        &        &        &        &        \\
  & \ddots &        &        &        &        \\
  &        & Q      &        &        &        \\
  &        &        & P      &        &        \\
  &        &        &        & R      &        \\
  &        &        &        &        & \ddots \\
  &        &        &        &        &        & R
\end{bmatrix}.
\]
So the quadratic-cost CFTOC problem can be written as a standard QP in which the dynamics
appear as equality constraints, the state and input limits appear as inequalities, and the cost is
a quadratic function of the stacked state-input vector \(z\). This formulation stays very close to
the original optimal control problem, since both predicted states and predicted inputs remain
explicit throughout.
\subsection{QP formulation by substitution}

A second way to rewrite the quadratic-cost CFTOC problem as a quadratic program is to
\emph{substitute the state equations} directly into the cost and the constraints, so that the only
decision variables are the future inputs. In this formulation, one starts from
\[
x_{i+1}=Ax_i+Bu_i,\qquad x_0=x(k),
\]
and eliminates the predicted states \(x_1,\dots,x_N\) recursively.\\
Define the stacked input vector
\[
U=
\begin{bmatrix}
u_0^\top & \cdots & u_{N-1}^\top
\end{bmatrix}^\top.
\]

\paragraph{Step 1: Rewrite the cost.}

After substituting the predicted states into the quadratic objective, the cost can be written as
\[
J(x(k),U)=U^\top H U + 2x(k)^\top F U + x(k)^\top Y x(k).
\]
Equivalently,
\[
J(x(k),U)=
\begin{bmatrix}
U^\top & x(k)^\top
\end{bmatrix}
\begin{bmatrix}
H & F^\top\\
F & Y
\end{bmatrix}
\begin{bmatrix}
U\\
x(k)
\end{bmatrix}.
\]
Since the stage and terminal costs are assumed nonnegative, this block matrix satisfies
\[
\begin{bmatrix}
H & F^\top\\
F & Y
\end{bmatrix}\succeq 0.
\]

\paragraph{Step 2: Rewrite the constraints compactly.}

The input, state, and terminal constraints can also be expressed directly in terms of \(U\) and
the current state \(x(k)\). This gives the compact inequality description
\[
GU \le w + Ex(k).
\]
To make this more explicit, suppose that
\[
\cX=\{x\mid A_xx\le b_x\},\qquad
\cU=\{u\mid A_uu\le b_u\},\qquad
\cX_f=\{x\mid A_fx\le b_f\}.
\]
Then the matrix \(G\), the vector \(E\), and the right-hand side \(w\) are obtained by stacking
the input constraints first, then the state constraints along the horizon, and finally the terminal
constraint. In block form,
\[
G=
\begin{bmatrix}
A_u & 0 & \cdots & 0\\
0 & A_u & \cdots & 0\\
\vdots & \vdots & \ddots & \vdots\\
0 & 0 & \cdots & A_u\\
0 & 0 & \cdots & 0\\
A_xB & 0 & \cdots & 0\\
A_xAB & A_xB & \cdots & 0\\
\vdots & \vdots & \ddots & \vdots\\
A_fA^{N-1}B & A_fA^{N-2}B & \cdots & A_fB
\end{bmatrix},
\]
\[
E=
\begin{bmatrix}
0\\
0\\
\vdots\\
0\\
-A_x\\
-A_xA\\
-A_xA^2\\
\vdots\\
-A_fA^N
\end{bmatrix},
\qquad
w=
\begin{bmatrix}
b_u\\
b_u\\
\vdots\\
b_u\\
b_x\\
b_x\\
b_x\\
\vdots\\
b_f
\end{bmatrix}.
\]
This compact form comes from writing each predicted state as an affine function of \(x(k)\) and
the input sequence \(U\), and then imposing the polyhedral constraints on those expressions.

\paragraph{Step 3: Rewrite the optimal control problem as a QP.}

The quadratic-cost CFTOC problem therefore becomes
\[
J^\star(x(k))
=
\min_U
\begin{bmatrix}
U^\top & x(k)^\top
\end{bmatrix}
\begin{bmatrix}
H & F^\top\\
F & Y
\end{bmatrix}
\begin{bmatrix}
U\\
x(k)
\end{bmatrix}
\]
subject to
\[
GU \le w + Ex(k).
\]
This is a quadratic program in the variable \(U\), parameterized by the current state \(x(k)\).
For a given \(x(k)\), the optimal input sequence \(U^\star\) can therefore be obtained by solving
this QP.

\subsection{Quadratic-cost state-feedback solution}

Consider again the quadratic program obtained after substitution:
\[
J^\star(x(k))
=
\min_U
\begin{bmatrix}
U^\top & x(k)^\top
\end{bmatrix}
\begin{bmatrix}
H & F^\top\\
F & Y
\end{bmatrix}
\begin{bmatrix}
U\\
x(k)
\end{bmatrix}
\]
subject to
\[
GU \le w + Ex(k).
\]
For each fixed current state \(x(k)\), this problem can be solved numerically and yields an
optimal input sequence
\[
U^\star(x(k))
=
\begin{bmatrix}
u_0^\star(x(k))^\top & \cdots & u_{N-1}^\star(x(k))^\top
\end{bmatrix}^\top.
\]
As in MPC only the first component of the optimal sequence is applied to the plant, the
relevant feedback law is
\[
u_0^\star=\kappa(x(k)).
\]
The important point is that this QP is not just an optimization problem with state \(x(k)\)
inserted into it numerically. It is a \emph{multiparametric quadratic program} (mp-QP), where
the current state plays the role of a parameter. Because of this, the optimizer has a very
specific structure as a function of the state.\\

In particular, on the feasible set \(\mathcal X_0\), the first component of the optimal solution
defines a map
\[
\kappa:\R^n\to\R^m,
\]
which is continuous and piecewise affine on polyhedra. More precisely, there exists a finite
collection of polyhedral regions
\[
CR^j=\{x\in\R^n \mid H^j x \le K^j\},\qquad j=1,\dots,N^r,
\]
which form a partition of the feasible polyhedron \(\mathcal X_0\), such that on each region one
has
\[
\kappa(x)=F^j x + g^j,\qquad x\in CR^j,\quad j=1,\dots,N^r.
\]
So even though the original optimal control problem is constrained and does not admit a single
global linear feedback law, it still has a highly structured solution: the state space is divided
into polyhedral regions, and on each region the optimal controller is affine.\\

The optimal value function has a similarly regular form. The map \(J^\star(x(k))\), which
assigns to each feasible initial state the optimal cost of the finite-horizon problem, is convex
and piecewise quadratic on polyhedra. Thus, both the control law and the value function inherit
a polyhedral structure from the parametric QP formulation.\\
This observation is the starting point for \emph{explicit MPC}. Instead of solving the QP online
at every sampling instant, one may compute offline the polyhedral partition together with the
corresponding affine control laws. The online step is then reduced to locating the current state
in the appropriate region and evaluating the affine formula associated with that region.

\begin{theorembox}[Quadratic-cost state-feedback structure]
For quadratic-cost CFTOC, the parametric optimization problem is an mp-QP. As a result, the
first optimal input defines a continuous piecewise affine state-feedback law on the feasible set,
while the associated optimal value function is convex and piecewise quadratic.
\end{theorembox}
\section{CFTOC with \(1\)-norm and \(\infty\)-norm costs}

So far we have considered quadratic-cost CFTOC, which leads naturally to a quadratic
program. One can treat alternative cost functions in a very similar
predictive-control framework. In particular, if the stage and terminal penalties are based on the
\(1\)-norm or the \(\infty\)-norm, the constrained finite time optimal control problem can be
rewritten as a linear program.\\
More precisely, consider the finite-horizon problem
\[
J^\star(x(k))
=
\min_U
\norm{Px_N}_p+\sum_{i=0}^{N-1}\left(\norm{Qx_i}_p+\norm{Ru_i}_p\right),
\qquad p\in\{1,\infty\},
\]
subject to
\[
x_{i+1}=Ax_i+Bu_i,\qquad i=0,\dots,N-1,
\]
\[
x_i\in\mathcal X,\qquad u_i\in\mathcal U,\qquad i=0,\dots,N-1,
\]
\[
x_N\in\mathcal X_f,
\qquad
x_0=x(k).
\]
The key point is that both the \(1\)-norm and the \(\infty\)-norm admit an equivalent
\emph{epigraph formulation}, in which the norm is replaced by a linear objective together with
additional linear inequalities. This is what makes it possible to transform the CFTOC problem
into an LP.

\subsection{\(\ell_\infty\)-minimization}

To see the basic idea, it is useful to start from the simpler constrained minimization problem
\[
\min_{x\in\R^n}\norm{x}_\infty
\qquad \text{subject to}\qquad
Fx\le g.
\]
Since
\[
\norm{x}_\infty=\max\{|x_1|,\dots,|x_n|\},
\]
this problem can be viewed as the minimization of the maximum of finitely many linear
functions:
\[
\min_{x\in\R^n}
\max\{x_1,\dots,x_n,-x_1,\dots,-x_n\}
\qquad \text{subject to}\qquad
Fx\le g.
\]
The standard trick is then to introduce a scalar auxiliary variable \(t\) that upper-bounds all
components in absolute value. This gives the equivalent problem
\[
\min_{x,t} t
\]
subject to
\[
x_i\le t,\qquad i=1,\dots,n,
\]
\[
-x_i\le t,\qquad i=1,\dots,n,
\]
\[
Fx\le g.
\]
In compact vector form, this can be written as
\[
\min_{x,t} t
\]
subject to
\[
-\mathbf 1\, t \le x \le \mathbf 1\, t,
\qquad
Fx\le g.
\]
Here \(\mathbf 1\) denotes the vector of ones. The constraint
\[
-\mathbf 1\, t \le x \le \mathbf 1\, t
\]
simply says that every component of \(x\) must satisfy \(|x_i|\le t\), so the scalar \(t\)
bounds the absolute value of all entries simultaneously. Minimizing \(t\) is therefore exactly
the same as minimizing \(\norm{x}_\infty\).\\
The \(\infty\)-norm is replaced by a linear cost in an
additional scalar variable together with linear inequalities. The problem is no longer written in
terms of a norm objective, but it has become a linear program.

\subsection{\(\ell_1\)-minimization}

The same idea works for the \(1\)-norm. Consider
\[
\min_{x\in\R^n}\norm{x}_1
\qquad \text{subject to}\qquad
Fx\le g.
\]
Since
\[
\norm{x}_1=\sum_{i=1}^n |x_i|,
\]
one may write this as
\[
\min_{x\in\R^n}
\sum_{i=1}^n \max\{x_i,-x_i\}
\qquad \text{subject to}\qquad
Fx\le g.
\]
Now, instead of a single scalar epigraph variable, one introduces a vector
\[
t\in\R^n
\]
whose components bound the absolute value of the corresponding entries of \(x\). The problem
becomes
\[
\min_{x,t} t_1+\cdots+t_n
\]
subject to
\[
x_i\le t_i,\qquad i=1,\dots,n,
\]
\[
-x_i\le t_i,\qquad i=1,\dots,n,
\]
\[
Fx\le g.
\]
In compact form,
\[
\min_{x,t} \mathbf 1^\top t
\]
subject to
\[
-t\le x\le t,
\qquad
Fx\le g.
\]
Again, \(\mathbf 1\) denotes the vector of ones. The constraint
\[
-t\le x\le t
\]
means that each component \(t_i\) bounds the absolute value of the corresponding component
\(x_i\). At the optimum one therefore has \(t_i=|x_i|\), and minimizing \(\mathbf 1^\top t\)
is equivalent to minimizing \(\norm{x}_1\).\\
So in both cases the same mechanism is at work: one introduces additional variables and writes
the norm minimization problem in epigraph form. This is exactly the step that later allows the
finite-horizon control problem to be expressed as an LP.

\subsection{Construction of the LP with substitution}

Returning now to CFTOC, we can proceed exactly as in the quadratic case with
substitution: one eliminates the predicted states and writes everything directly in terms of the
current state \(x(k)\), the future inputs, and the auxiliary epigraph variables.\\
For the \(\infty\)-norm case, one introduces scalar variables
\[
\varepsilon_i^x,\qquad i=0,\dots,N,
\]
and
\[
\varepsilon_i^u,\qquad i=0,\dots,N-1,
\]
which bound the norms of the weighted state and input terms. Using the substituted prediction
formula for the states, the problem can be written equivalently as
\[
\min_z
\varepsilon_0^x+\cdots+\varepsilon_N^x+\varepsilon_0^u+\cdots+\varepsilon_{N-1}^u
\]
subject to inequalities of the form
\[
-\mathbf 1_n\,\varepsilon_i^x
\le
\pm Q\!\left(A^i x_0+\sum_{j=0}^{i-1}A^jBu_{i-1-j}\right),
\]
\[
-\mathbf 1_r\,\varepsilon_N^x
\le
\pm P\!\left(A^N x_0+\sum_{j=0}^{N-1}A^jBu_{N-1-j}\right),
\]
\[
-\mathbf 1_m\,\varepsilon_i^u \le \pm Ru_i,
\]
together with the state, input, and terminal constraints
\[
A^i x_0+\sum_{j=0}^{i-1}A^jBu_{i-1-j}\in\mathcal X,
\qquad
u_i\in\mathcal U,
\]
\[
A^N x_0+\sum_{j=0}^{N-1}A^jBu_{N-1-j}\in\mathcal X_f,
\qquad
x_0=x(k).
\]
The detailed block matrices are not especially illuminating in prose, but the important message
is that after substitution and introduction of the epigraph variables, the entire problem has only
a linear objective and linear constraints.\\
This yields the standard LP form
\[
\min_z c^\top z
\qquad \text{subject to}\qquad
\bar G z \le \bar w + \bar S x(k),
\]
where the decision vector contains both the auxiliary variables and the input sequence:
\[
z
:=
\{\varepsilon_0^x,\dots,\varepsilon_N^x,\varepsilon_0^u,\dots,\varepsilon_{N-1}^u,
u_0^\top,\dots,u_{N-1}^\top\}\in\R^s.
\]
The matrices can then be partitioned :
\[
\bar G=
\begin{bmatrix}
G_\varepsilon & G_u\\
0 & G
\end{bmatrix},
\qquad
\bar S=
\begin{bmatrix}
S_\varepsilon\\
S
\end{bmatrix},
\qquad
\bar w=
\begin{bmatrix}
w_\varepsilon\\
w
\end{bmatrix}.
\]
For a given current state \(x(k)\), the optimal sequence \(U^\star\) can therefore be obtained
by solving a linear program. The \(1\)-norm case works in exactly the same spirit; only the
epigraph variables are vector-valued instead of scalar.

\subsection{\(1\)-/\(\infty\)-norm state-feedback solution}

Once the norm-based CFTOC problem has been written in the compact LP form
\[
\min_z c^\top z
\qquad \text{subject to}\qquad
\bar G z \le \bar w + \bar S x(k),
\]
the parametric structure of the optimizer becomes analogous to the quadratic case. The current
state \(x(k)\) again enters as a parameter, so the problem is a \emph{multiparametric linear
program} (mp-LP).\\
As a consequence, the first component of the multiparametric solution defines a state-feedback
law
\[
u_0^\star=\kappa(x(k)),\qquad \forall x(k)\in\mathcal X_0,
\]
where
\[
\kappa:\R^n\to\R^m
\]
is piecewise affine on polyhedra. More precisely, there exists a polyhedral partition
\[
CR^j=\{x\in\R^n\mid H^j x\le K^j\},\qquad j=1,\dots,N^r,
\]
of the feasible polyhedron \(\mathcal X_0\) such that
\[
\kappa(x)=F^j x + g^j,
\qquad
x\in CR^j,\quad j=1,\dots,N^r.
\]
If the optimizer is unique, this control law is continuous and piecewise affine. If the LP has
multiple optimizers, one can still construct a piecewise affine control law, but uniqueness is no
longer guaranteed automatically.\\
The value function also reflects the underlying linear structure. In contrast with the quadratic
case, where the value function is piecewise quadratic, in the LP case the optimal value
\[
J^\star(x(k))
\]
is convex and piecewise affine on polyhedra.

\subsection{Solution properties: quadratic versus \(1\)-/\(\infty\)-norm cost}

We conclude this part by contrasting the geometry of quadratic and linear costs. This
comparison is useful because it explains why the value functions and optimizer properties differ.\\

Let \(n\) denote the number of optimization variables. If the cost is quadratic and positive
definite, then the solution is always unique whenever the problem is feasible. More precisely,
the optimizer is either in the interior of the feasible set, in which case no inequality constraint
is active, or on the boundary of the feasible set, in which case at least one constraint is active.\\
For a linear cost, the situation is different. The problem may be unbounded. If it is bounded
and the optimum is unique, then it is attained at a vertex of the feasible set, which means that
at least \(n\) constraints are active. It may also happen that there is not a unique minimizer,
but rather a whole set of optimal solutions; in that case at least one constraint is active and the
optimum is attained along a face of the feasible polyhedron.\\

So the difference between the two formulations is not only algebraic but also geometric. A
strictly convex quadratic objective selects a unique point, while a linear objective may select an
entire face. This is why quadratic-cost CFTOC leads to a convex piecewise quadratic value
function and typically unique optimizers, whereas \(1\)-/\(\infty\)-norm CFTOC leads to a convex
piecewise affine value function and may admit multiple optimal solutions.

\begin{theorembox}[Norm-based CFTOC as an LP]
If the stage and terminal costs are based on the \(1\)-norm or the \(\infty\)-norm, the CFTOC
problem can be rewritten in epigraph form and transformed into a linear program. The resulting
parametric problem is an mp-LP, whose first optimal input is piecewise affine on a polyhedral
partition of the feasible set, while the value function is convex and piecewise affine.
\end{theorembox}

\section{Receding-horizon notation and control law}

Consider the discrete-time model
\[
x(k+1)=Ax(k)+Bu(k),
\qquad
y(k)=Cx(k),
\]
with
\[
x(k)\in\mathcal X,\qquad u(k)\in\mathcal U,\qquad \forall k\ge 0.
\]
Up to this point, we have mostly written the constrained finite time optimal control problem in
a compact form that suppresses the prediction origin. For MPC, however, it is useful to make
this dependence explicit, since the optimization problem is solved repeatedly at successive time
instants.

\subsection{CFTOC problem solved at time \(k\)}

At time \(k\), the constrained finite time optimal control problem is written as
\[
J_{k\to k+N|k}^\star(x(k))
=
\min_{U_{k\to k+N|k}}
\terminalcost(x_{k+N|k})
+
\sum_{i=0}^{N-1}
\stagecost(x_{k+i|k},u_{k+i|k})
\]
subject to
\[
x_{k+i+1|k}=Ax_{k+i|k}+Bu_{k+i|k},
\qquad i=0,\dots,N-1,
\]
\[
x_{k+i|k}\in\mathcal X,\qquad u_{k+i|k}\in\mathcal U,
\qquad i=0,\dots,N-1,
\]
\[
x_{k+N|k}\in\mathcal X_f,
\]
\[
x_{k|k}=x(k),
\]
where the decision sequence is
\[
U_{k\to k+N|k}
=
\{u_{k|k},\dots,u_{k+N-1|k}\}.
\]
So the problem is solved at time \(k\), starting from the currently measured or estimated state
\(x(k)\), and it produces an optimal input sequence over the horizon \(k,\dots,k+N-1\).

\subsection{Meaning of the prediction notation}

The notation is best understood carefully. The symbol \(x(k)\) denotes the actual state of the
system at time \(k\). By contrast,
\[
x_{k+i|k}
\]
denotes the state predicted for time \(k+i\), where the prediction is made at time \(k\). This
predicted state is obtained by starting from the current state
\[
x_{k|k}=x(k)
\]
and propagating the model forward using the planned inputs
\[
u_{k|k},\dots,u_{k+i-1|k}.
\]
For instance, the first predicted step is
\[
x_{k+1|k}=Ax_{k|k}+Bu_{k|k}.
\]
It is important to keep in mind that the same physical time instant may correspond to different
predictions, depending on when the prediction is made. For example, \(x_{3|1}\) is the state
predicted for time \(3\) using information available at time \(1\), whereas \(x_{3|2}\) is the state
predicted for time \(3\) using information available at time \(2\). In general, these two
quantities are different because they are based on different initial conditions and different input
plans.\\
Similarly,
\[
u_{k+i|k}
\]
is read as “the input at time \(k+i\), computed at time \(k\).” This notation keeps track of the
fact that in MPC one computes an entire future plan, even though not all of it will actually be
implemented.

\subsection{From the optimal sequence to the applied input}

Let
\[
U_{k\to k+N|k}^\star
=
\{u_{k|k}^\star,\dots,u_{k+N-1|k}^\star\}
\]
denote the optimal solution of the CFTOC problem solved at time \(k\). Although this optimal
sequence contains \(N\) future control moves, only its first element is applied to the plant:
\[
u(k)=u_{k|k}^\star(x(k)).
\]
This is the defining feature of receding-horizon control. The optimizer computes a whole
sequence, but the controller implements only the first input and discards the rest after one
sampling interval.

At the next time instant, the state is updated to
\[
x_{k+1|k+1}=x(k+1),
\]
and the optimization problem is solved again from this new initial condition. In this way, the
controller continuously replans as new state information becomes available.

\subsection{Receding-horizon control law and closed loop}

The first optimal move defines the receding-horizon feedback law
\[
\kappa_k(x(k)) = u_{k|k}^\star(x(k)).
\]
The corresponding closed-loop system is therefore
\[
x(k+1)=Ax(k)+B\kappa_k(x(k))=:g_{\mathrm{cl}}(x(k)),
\qquad k\ge 0.
\]
So even though the optimization problem is finite-horizon and open-loop in its internal
prediction, the repeated re-solution at each step turns the overall control scheme into a
feedback controller.

\subsection{Time-invariant systems}

If the system, the constraints, and the cost function are time-invariant, then the CFTOC
problem has the same structure at every time \(k\). In that case, the feedback law obtained by
receding-horizon optimization also becomes time-invariant, and the notation can be simplified.\\
One then writes
\[
J^\star(x(k))
=
\min_U
\terminalcost(x_N)
+
\sum_{i=0}^{N-1}\stagecost(x_i,u_i)
\]
subject to
\[
x_{i+1}=Ax_i+Bu_i,\qquad i=0,\dots,N-1,
\]
\[
x_i\in\mathcal X,\qquad u_i\in\mathcal U,\qquad i=0,\dots,N-1,
\]
\[
x_N\in\mathcal X_f,
\qquad
x_0=x(k),
\]
where
\[
U=\{u_0,\dots,u_{N-1}\}.
\]
In the same spirit, the control law is written simply as
\[
u(k)=\kappa(x(k))
\]
instead of \(\kappa_k(x(k))\).

\begin{theorembox}[Receding-horizon principle]
At each sampling instant, MPC solves a finite-horizon optimal control problem, applies only the
first element of the optimal input sequence, and then repeats the optimization at the next time
step using the new state. This repeated re-optimization is what gives MPC its feedback nature.
\end{theorembox}
\section{Summary}

Constrained finite time optimal control is the optimization problem at the heart of nominal
MPC. Starting from a linear discrete-time model, one predicts future trajectories over a finite
horizon, penalizes those trajectories through a chosen cost function, and imposes state, input,
and terminal constraints explicitly. The resulting feasible set consists of the initial states from
which at least one admissible predicted trajectory exists.\\

When the cost is quadratic and the constraints are polyhedral, the problem can be written as a
convex quadratic program. This may be done either by keeping the state variables explicitly and
enforcing the dynamics as equality constraints, or by substituting the dynamics and optimizing
only over the input sequence. In the parametric viewpoint, the first optimal input is a piecewise
affine function of the state and the value function is piecewise quadratic.\\

When the cost is based on the \(1\)-norm or the \(\infty\)-norm, auxiliary epigraph variables allow
the problem to be rewritten as a linear program. In that case the first optimal input is again
piecewise affine, but the value function becomes piecewise affine and multiple optimal solutions
may occur.\\

Finally, the finite-horizon problem is embedded into a closed-loop strategy through the
receding-horizon principle. At each sampling instant, one solves a CFTOC problem from the
current state, applies only the first input of the optimal sequence, and then repeats the same
procedure at the next time step. This is the basic operational mechanism of model predictive
control.